\documentclass[10pt]{article}

\usepackage[utf8]{inputenc}

\usepackage[T1]{fontenc}

\usepackage{amsmath}

\usepackage{amsfonts}

\usepackage{amssymb}

\usepackage[version=4]{mhchem}

\usepackage{stmaryrd}

\usepackage{graphicx}

\usepackage[export]{adjustbox}

\graphicspath{ {./images/} }



\begin{document}

ного знака. Существование У . н. а. доказал А. А. Марков (см. [1]). Важной характеристикой У. н. а. является его с л о ж п о т ь, т. е. длина его изображения (см. также Алгоритма сложность описания). Для минимальной сложности У. н. а. как функции от $n$ (количества символов в алфавите $A$ ) получены отличающиеся лишь на аддитивную константу нижняя и верхняя оценки вида $5 n+C$ (см. [2]).



\includegraphics[max width=\textwidth, center]{2023_07_09_351bf4e4e0df7a4c75cbg-1}

1954 (Тр. Матем. ин-та АН СССР, т 42); [2] ЖК а р о в В. Г., О рия алгорифмов и математическая логика, М., 1974 , с. $34-54$.



УнивеРСАлЫНІ РЯД С. Н. Артемов.



\[

\sum_{i=1}^{\infty} \varphi_{i}(x), \quad x \in[a, b]

\]



с помощью к-рого могут быть представлены в том или ином смысле все функции заданного класса. Напр., существует такой ряд (1), что для каждой непрерывной на $[a, b]$ функции $f$ найдется подпоследовательность частных сумм этого ряда $\sum_{i=1}^{n_{k}} \varphi_{i}(x)$, сходящаяся $\mathrm{\kappa} f(x)$ равномерно на $[a, b]$.



Существуют тригонометрические ряды



\[

\frac{a_{0}}{2}+\sum_{i=1}^{\infty}\left(a_{i} \cos i x+b_{i} \sin i x\right)

\]



со стремящимися к нулю коэффициентами такие, что для каждой измеримой (по Лебегу) на $[0,2 \pi]$ функции $f$ имеется подпоследовательность частных сумм ряда (2), сходящаяся $\kappa ~ f(x)$ почти всюду.



Указанные ряды наз. универсальными относительно подпоследовательностей частных сумм. Рассматриваются также др. определения У. р. Напр., ряды (1), универсальные относительно подрядов $\sum_{k=1}^{\infty} \varphi_{i_{k}}(x)$ или относительно перестановок членов ряда (1).



Лum. [1] А л е к с и ч Г., Проблемы сходимости ортогональных рядов, пер. с англ., М., 1963 , [2] Б а р и Н. К., Тригонометрические ряды, М., 1961, [3] Т а ЈІа л я н А. А., «Успехи матем. наук», 1960 , т. 15 , № 5, с $77 \frac{141}{\text { C. }}$.



УНИВЕРСАЛЬНЫХ АЛГЕБР МНОГООБРАВИЕ класс универсальных алгебр, определяемый системой тождеств (ср. Алгебраических систем многообразие). У. а. м. характеризуется как непустой класс алгебр, замкнутый относительно факторалгебр, подалгебр и прямых произведений. Последние два условия можно заменить требованием замкнутости относительно подпрямых произведений. У. а. м. наз. т р и в и а л ьны м, если оно состоит из одноәлементных алгебр. Каждое нетривиальное У . а. м. содержит свободную алгебру с базой любой мощности. Если $X$ и $Y$ - базы одной и той же свободной алгебры нетривиального У . а. м. и $X$ бесконечна, то $X$ и $Y$ равномопны. Требование бесконечности одной из баз существенно. Оно может быть снято, если У. а. м. содержит неодноәлементную конечную алгебру.



У. а. м., порожденное классом $K$, состоит из всех факторалгебр всевозможных подпрямых произведений алгебр из $K$. Все конечно порожденные алгебры из У. а. м., порожденного конечной алгеброй, конечны. Конгруәнции любой алгебры из У. а. м. $M$ сигнатуры $\Omega$ перестановочны в том и только в том случае, когда существует такой тернарный терм $f$ сигпатуры $\Omega$, что тождества



\[

f(x, x, y)=y=f(y, x, x)

\]



справедливы во всех алгебрах из $M$. Аналогичным образом могут быть охарактеризованы У. а. м., чьи алгебры обладают модулярными или дистрибутивными решетками конгруэнций (см. [1-4,7,9,10]).



В многообразии $M$ n-арная операция $f$ паз. т р ив и а ј ь н о й, если во всех алгебрах из $M$ справедливо тождество $f\left(x_{1}, . ., x_{n}\right)=f\left(y_{1}, \ldots, y_{n}\right)$. Напр., в многообразии колец с нулевым умножением операция умножения тривиальна. Каждую тривиальную операцию $f$ можно заменить 0 -арной операцией $v_{f}$, определяемой равенством $v_{f}=f\left(x_{1}, . ., x_{n}\right)$. Пусть сигнатуры $\Omega$ и $\Omega^{\prime}$ У. а. м. $M$ и $M^{\prime}$ соответственно не содержат тривиальных операций. Отображение $\Phi$ сигнатуры $\Omega$ в множество $W\left(\Omega^{\prime}\right)$ термов сигнатуры $\Omega^{\prime}$ наз. д о п у с т им ы м, если арности операций $f$ и Ф $(f)$ совпадают для всех $f \in \Omega$. Допустимое отображение $\Phi$ естественным образом шродолжается до отображения $W(\Omega)$ в $W\left(\Omega^{\prime}\right)$, к-рое также обозначается $\Phi$. Многообразия $M$ и $M^{\prime}$ наз. р а ци о на ль но-э к в и в а л е н т ным и, если существуют допустимые отображения $\Phi: \Omega \rightarrow$ $\rightarrow W\left(\Omega^{\prime}\right)$ и $\Phi^{\prime}: \Omega^{\prime} \rightarrow W(\Omega)$ такие, что $f=\Phi^{\prime}(\Phi(f))$ для всех $f \in \Omega, f^{\prime}=\Phi\left(\Phi^{\prime}\left(f^{\prime}\right)\right)$ для всех $f^{\prime} \in \Omega^{\prime}$ и для каждого тождества $u=v$ (соответственно $u^{\prime}=v^{\prime}$ ), входящего в определение многообразия $M$ (многообразия $M^{\prime}$ ), тождество $\Phi(u)=\Phi(v)$ (соответственно $\Phi^{\prime}\left(u^{\prime}\right)=\Phi^{\prime}\left(v^{\prime}\right)$ ) сграведливо во всех алгебрах из $M^{\prime}$ (из $M$ ). Последнее требование равносильно тому, что каждая алгебра $A$ из $M\left(A^{\prime}\right.$ из $\left.M^{\prime}\right)$ превращается в алгебру из $M^{\prime}$ (из $\left.M\right)$, если каждую $n$-арную операцию $f^{\prime}$ из $\Omega^{\prime}(f$ из $\Omega$ ) определить равенством $f^{\prime}\left(x_{1}, \ldots, x_{n}\right)=\Phi^{\prime}\left(f^{\prime}\right)\left(x_{1}, \ldots, x_{n}\right)$ (соответственно $\left.f\left(x_{1}, \ldots, x_{n}\right)=\Phi(f)\left(x_{1}, \ldots, x_{n}\right)\right)$. $\mathrm{Pa}-$ ционально эквивалентны многообразия булевых колец и булевых алгеб $p$. Многообразше унар.ыи алгеб $p$ сигнатуры $\Omega$, определяемое тождествами



\[

\left\{u_{\iota}(x)=v_{\iota}(x) \mid \iota \in \Im\right\}

\]



рационально эквивалентно многообразию всех левых $R$-полигонов, где $R$ - фактормоноид свободного моноида со свободной порождающей системой $\Omega$ по конгруәнции, порожденной всеми парами $\left\{\left(u_{\iota}, v_{\iota}\right) \mid \iota \in \Im\right.$ $\}$. У. а. м. $M$ рационально эквивалентно многообразию всех правых модулей над нек-рым ассоциативным кольцом тогда и только тогда, когда конгруэнции любой алгебры из $M$ перестановочны, конечные свободные произведения в $M$ совпадают с прямыми произведениями п существует нуль-арная производная операция, отмечающая подалгебру. Первые два условия можно заменить требованием: каждая подалгебра любой алгебры из $M$ является классом нек-рой конгруэнции и каждая конгруэнция любой алгебры из $M$ однозначно определяется своим классом, являющимся подалгеброй $[3,5,6,7]$.



Многообразие репеток, порожденное решетками конгруәнций всех алгебр нек-рого У. а. м., наз. к о н г р уэ н ц-м н о г о о бр а з и е м. Не всякое многообразие решеток является конгруэнц-многообразием. Существуют не модулярные конгруэнц-многообразия, отличные от многообразия всех решеток [7, 8].



Лит [1] K о н П., Универсальная алгебра, пер. с англ.. M., 1968, [2] к у p о ш А. Г., Јекции по общей алгебре, 2 изд., М $1970,[4]$, С к о р н я к о в Ј. А., Элементы общей алгебры, М., 1983; L5] Ч а к a ғ b Б., "Acta Scient. Math.", 1963, t. 24, Nㅑ $1-2$,



\includegraphics[max width=\textwidth, center]{2023_07_09_351bf4e4e0df7a4c75cbg-1(1)}

B., 1979; [8] J ón s s o n B., "Coll. Math. Soc. J., Bolyai", 1982 , v. 29 , p. 421 - 36, [9] S m i t h J. D. H., Mal'cev varieties, B.- Hdlb.- N. Y., 1976; [10] T a y l o r W., "Algebra univer-

salis», 1973, v. 3, №, 3, p. 351-97.



УНИКУРСАЛЬНАЯ КРИВАЯ - плоская кривая $\Gamma$, к-рую можно обойти, побывав дважды только в точках самопересечения. Для того чтобы кривая была уникурсальной, необходимо и достаточно, чтобы у нее было не более двух точек, через к-рые проходит нечетное число путей. Если $\Gamma$ - плоская алгебрапч. кривая $n$-го порядка, имеющая максимальное число $\delta$ двойных точек (включая несобственные и мнимые), то $\delta=$ $=\frac{(n-1)(n-2)}{2}$ (причем точки кратности $k$ рассматриваются как $\frac{k(k-1)}{2}$ двойных точек).





\end{document}